\subsection{Identitas dan Invers}

Ada sebuah matriks khusus, yang dinyatakan dengan $I$, dan disebut
\index{matriks!identitas} \textbf{matriks identitas}. Matriks identitas selalu merupakan matriks
persegi, dan memiliki
sifat bahwa diagonal utamanya berisi angka satu dan semua posisi lainnya
berisi nol. Berikut beberapa matriks identitas dengan berbagai ukuran.
\begin{equation*}
\leftB 1\rightB ,\leftB
\begin{array}{cc}
1 & 0 \\
0 & 1
\end{array}
\rightB ,\leftB
\begin{array}{ccc}
1 & 0 & 0 \\
0 & 1 & 0 \\
0 & 0 & 1
\end{array}
\rightB ,\leftB
\begin{array}{cccc}
1 & 0 & 0 & 0 \\
0 & 1 & 0 & 0 \\
0 & 0 & 1 & 0 \\
0 & 0 & 0 & 1
\end{array}
\rightB 
\end{equation*}
Matriks pertama adalah matriks identitas $1\times 1$, matriks kedua adalah matriks identitas $2\times 2$,
dan seterusnya. Dengan memperluas pola ini, Anda mungkin dapat melihat
bentuk matriks identitas $n\times n$. Jika perlu dibedakan
ukuran matriks identitas yang sedang dibahas, kita akan menggunakan
notasi $I_n$ untuk matriks identitas $n \times n$. 

Matriks identitas begitu penting sehingga terdapat
simbol khusus untuk menyatakan entri pada posisi $(i,j)$ dari matriks identitas. Simbol ini diberikan oleh
$I_{ij}=\delta _{ij}$, dengan $\delta _{ij}$ sebagai \textbf{simbol Kronecker}
yang didefinisikan
\index{simbol Kronecker} oleh
\begin{equation*}
\delta _{ij}=\left\{
\begin{array}{c}
1
\text{ jika }i=j \\
0\text{ jika }i\neq j
\end{array}
\right.
\end{equation*}

$I_n$ disebut \textbf{matriks identitas} karena merupakan \textbf{identitas perkalian} dalam pengertian berikut.

\begin{lemma}{Perkalian dengan Matriks Identitas}\label{lem:multbyidentity}
Misalkan $A$ adalah matriks $m\times n$ dan $I_{n}$ adalah matriks identitas
$n\times n$\textbf{.} Maka $AI_{n}=A.$ Jika $I_{m}$ adalah matriks identitas
$m\times m$, berlaku pula $I_{m}A=A.$
\end{lemma}

\ifdefined\showproofs
\begin{proof}
Entri $(i,j)$ dari $AI_n$ diberikan oleh:
\begin{equation*}
\sum_{k}a_{ik}\delta _{kj}=a_{ij}
\end{equation*}
sehingga $AI_{n}=A.$ Kasus lainnya diserahkan sebagai latihan untuk Anda. 
\end{proof}
\fi
Sekarang kita mendefinisikan operasi matriks yang dalam beberapa hal berperan seperti pembagian. 

\begin{definition}{Invers Suatu Matriks}\label{def:invertiblematrix}
Matriks persegi $n\times n$, yaitu $A$, dikatakan memiliki \textbf{invers} $A^{-1}$
jika dan hanya jika
\begin{equation*}
AA^{-1}=A^{-1}A=I_n
\end{equation*}
Dalam hal ini, matriks $A$ disebut
\index{matriks!invers} \textbf{invertibel}.
\index{matriks!invertibel}
\end{definition}

Matriks $A^{-1}$ tersebut memiliki ukuran yang sama dengan matriks $A$.
Sangat penting untuk diperhatikan bahwa invers suatu matriks, jika ada,
bersifat unik. Cara lain untuk memahaminya adalah: jika suatu matriks bertindak seperti invers,
maka matriks itu $\textbf{memang}$ inversnya.

\begin{theorem}{Ketunggalan Invers}\label{thm:uniqueinverse}
Misalkan $A$ adalah matriks $n \times\ n$ sedemikian sehingga invers $A^{-1}$ ada. Maka hanya ada satu
matriks invers semacam itu.
Artinya, untuk setiap matriks $B$ sedemikian sehingga $AB=BA=I$, berlaku $B=A^{-1}$.
\end{theorem}

\ifdefined\showproofs
\begin{proof} Dalam bukti ini, diasumsikan bahwa $I$ adalah matriks identitas $n \times n$.
Misalkan $A, B$ adalah matriks $n \times n$ sedemikian sehingga $A^{-1}$ ada dan $AB=BA=I$.
Kita ingin menunjukkan bahwa $A^{-1} = B$.
Sekarang, dengan menggunakan sifat-sifat yang telah kita lihat, diperoleh:
\begin{equation*}
A^{-1}=A^{-1}I=A^{-1}\left( AB\right) =\left( A^{-1}A\right) B=IB=B
\end{equation*}
Jadi, $A^{-1} = B$, yang menunjukkan bahwa invers tersebut tunggal.
\end{proof}
\fi
Contoh berikut menunjukkan cara memeriksa invers suatu matriks. 

\begin{example}{Memeriksa Invers Suatu Matriks}\label{exa:verifyinginverse}
Misalkan $A=\leftB
\begin{array}{rr}
1 & 1 \\
1 & 2
\end{array}
\rightB .$ Tunjukkan bahwa $\leftB
\begin{array}{rr}
2 & -1 \\
-1 & 1
\end{array}
\rightB $ adalah invers dari $A.$
\end{example}

\begin{solution} Untuk memeriksanya, kalikan
\begin{equation*}
\leftB
\begin{array}{rr}
1 & 1 \\
1 & 2
\end{array}
\rightB \leftB
\begin{array}{rr}
2 & -1 \\
-1 & 1
\end{array}
\rightB =\allowbreak \leftB
\begin{array}{rr}
1 & 0 \\
0 & 1
\end{array}
\rightB = I 
\end{equation*}
dan
\begin{equation*}
\leftB
\begin{array}{rr}
2 & -1 \\
-1 & 1
\end{array}
\rightB \leftB
\begin{array}{rr}
1 & 1 \\
1 & 2
\end{array}
\rightB =\allowbreak \leftB
\begin{array}{rr}
1 & 0 \\
0 & 1
\end{array}
\rightB = I 
\end{equation*}
yang menunjukkan bahwa matriks ini memang invers dari $A.$
\end{solution}

Berbeda dengan perkalian bilangan biasa, dapat terjadi bahwa $A\neq 0$, tetapi
$A$ tidak memiliki invers. Hal ini diilustrasikan dalam contoh berikut.

\begin{example}{Matriks Taknol yang Tidak Memiliki Invers}\label{exa:noninvertiblematrix}
Misalkan $A=\leftB
\begin{array}{rr}
1 & 1 \\
1 & 1
\end{array}
\rightB .$ Tunjukkan bahwa $A$ tidak memiliki invers.
\end{example}

\begin{solution} Kita mungkin mengira $A$ memiliki invers karena tidak sama dengan nol.
Namun, perhatikan bahwa
\begin{equation*}
\leftB
\begin{array}{rr}
1 & 1 \\
1 & 1
\end{array}
\rightB \leftB
\begin{array}{r}
-1 \\
1
\end{array}
\rightB =\leftB
\begin{array}{r}
0 \\
0
\end{array}
\rightB
\end{equation*}
Jika $A^{-1}$ ada, kita akan memperoleh
\begin{eqnarray*}
\leftB
\begin{array}{r}
0 \\
0
\end{array}
\rightB &=& A^{-1}\left( \leftB
\begin{array}{r}
0 \\
0
\end{array}
\rightB \right) \\
&=& A^{-1}\left( A\leftB
\begin{array}{r}
-1 \\
1
\end{array}
\rightB \right) \\
&=&\left( A^{-1}A\right) \leftB
\begin{array}{r}
-1 \\
1
\end{array}
\rightB \\
&=&I\leftB
\begin{array}{r}
-1 \\
1
\end{array}
\rightB \\
&=&\leftB
\begin{array}{r}
-1 \\
1
\end{array}
\rightB
\end{eqnarray*}
Ini menyatakan bahwa
\begin{equation*}
\leftB
\begin{array}{r}
0 \\
0
\end{array}
\rightB
=
\leftB
\begin{array}{r}
-1 \\
1
\end{array}
\rightB
\end{equation*}
yang mustahil! Oleh karena itu, $A$ tidak memiliki invers. 
\end{solution}

Pada bagian berikutnya, kita akan membahas cara mencari invers suatu matriks, jika invers tersebut ada. 
